\chapter{Lectures}

Quelques livres qui m'ont servi, rangés selon les parties du livre. Aucun ne parle de
mathématiques assistées ; tous parlent d'un des métiers qu'on y croise.

\section*{Le métier de programmeur}
\begin{description}
  \item[Andy Hunt, David Thomas, \emph{The Pragmatic Programmer} (1999).] Le livre d'où
        viennent la plupart des habitudes des chapitres 1 et 3 : automatiser ce qu'on refait,
        garder ses sources en texte brut, ne pas laisser de fenêtre cassée.
  \item[Steve McConnell, \emph{Code Complete} (1993).] La construction du logiciel dans le
        détail : nommer, découper, relire. Valable mot pour mot pour un fichier Lean.
  \item[Martin Fowler, \emph{Refactoring} (1999).] Changer la forme sans changer le sens, par
        petits pas vérifiés. C'est la règle « un commit par idée » appliquée au code.
  \item[Michael Feathers, \emph{Working Effectively with Legacy Code} (2004).] Comment
        reprendre un ensemble qu'on n'a pas écrit, en posant d'abord des tests ; utile devant
        un dépôt qu'un assistant a fait grossir vite.
  \item[Chris Riccomini, Dmitriy Ryaboy, \emph{The Missing README} (2021).] Ce qu'on attend
        d'un nouvel ingénieur : pull requests, code review, astreintes. Une bonne
        introduction pour qui vient des mathématiques.
  \item[Rúnar Bjarnason, Paul Chiusano, \emph{Functional Programming in Scala} (2013).] La
        programmation fonctionnelle comme pont entre le code et la preuve.
  \item[Eric S. Raymond, \emph{The Art of UNIX Programming} (2003) ; Jeroen Janssens,
        \emph{Data Science at the Command Line} (2014).] Le terminal, les petits outils qu'on
        enchaîne, le texte comme format universel.
\end{description}

\section*{Faire tourner, vérifier, livrer}
\begin{description}
  \item[Gene Kim et al., \emph{The DevOps Handbook} (2016) ; Betsy Beyer et al.,
        \emph{Site Reliability Engineering} (2016).] L'intégration continue et le déploiement
        vus par ceux qui les ont industrialisés. On n'en garde qu'une petite partie, mais
        c'est la bonne.
  \item[Craig Silverman (dir.), \emph{Verification Handbook} (2014).] Écrit pour des
        journalistes : comment vérifier une source, une image, une citation avant de la
        publier. Les réflexes du chapitre~\ref{ch:litterature} y sont presque tous.
  \item[Jake Knapp, \emph{Sprint} (2016).] Une méthode en cinq jours pour tester une idée
        avant de la construire ; on en retient le découpage du temps et la décision écrite.
\end{description}

\section*{Travailler avec des modèles de langage}
\begin{description}
  \item[Chip Huyen, \emph{AI Engineering} (2025).] Construire avec des modèles de fondation :
        évaluation, contexte, agents. Le meilleur panorama technique que je connaisse.
\end{description}

\section*{Lire et écrire}
\begin{description}
  \item[Almuth Grésillon, \emph{Éléments de critique génétique. Lire les manuscrits
        modernes}.] Comment lire un manuscrit de travail, ses ratures et ses campagnes
        d'écriture. Le meilleur guide pour qui transcrit des notes, de mathématiques ou
        d'autre chose.
  \item[Raymond Queneau, \emph{Exercices de style} (1947).] Une même histoire en quatre-vingt-
        dix-neuf registres : la meilleure leçon sur ce qu'une passe d'humanisation change et
        ne doit pas changer.
  \item[Jonathan Schwabish, \emph{Better Data Visualizations} (2021) ; Donald Norman,
        \emph{The Design of Everyday Things} (1988).] Pour les figures des articles et les
        pages de l'interface locale.
\end{description}

\section*{Le métier de mathématicien}
\begin{description}
  \item[David Bessis, \emph{Mathematica} (2022) ; Pierre-Louis Lions, \emph{Dans la tête
        d'un mathématicien} (2020) ; Philippe Pajot, \emph{Parcours de mathématiciens}
        (2015).] Ce que fait un mathématicien quand il pense, dit par des mathématiciens.
        Aucun outil ne remplace ce travail ; les outils de ce livre servent à lui laisser
        le temps.
  \item[Roger Godement, \emph{Analyse mathématique} (1998--2001).] Un cours qui ne cache
        rien de ses raisons, et un modèle de prose mathématique à la première personne.
\end{description}
